% Generated draft: requires mathematical and source review.
\documentclass[11pt]{article}
\usepackage[margin=25mm]{geometry}
\usepackage{fontspec}
\setmainfont{Latin Modern Roman}
\usepackage{amsmath,amssymb,mathtools,unicode-math}
\setmathfont{Latin Modern Math}
\usepackage{enumitem,needspace,xurl,hyperref}
\setlength{\parindent}{0pt}
\setlength{\parskip}{5pt}
\setlength{\emergencystretch}{3em}
\newcommand{\N}{\mathbb{N}}
\newcommand{\Z}{\mathbb{Z}}
\newcommand{\Q}{\mathbb{Q}}
\newcommand{\R}{\mathbb{R}}
\newcommand{\C}{\mathbb{C}}
\newcommand{\F}{\mathbb{F}}
\newcommand{\A}{\mathbb{A}}
\newcommand{\PP}{\mathbb{P}}
\newcommand{\E}{\mathbb{E}}
\newcommand{\eps}{\varepsilon}
\newcommand{\dd}{\,\mathrm{d}}
\begin{document}
\section*{UnsolvedMath ID: 700637}
\textbf{Source problem code:} catalogue entry
\subsection{Definitions and mathematical statement}
Let \(G\) be a Chevalley group over a finite field \(\F_q\), where \(q\leq 9\), and suppose that its rank is at most \(6\). Let \(U,V\leq G\) be opposite unipotent subgroups, as specified in the source. For a subset \(S\subseteq G\), write \(\langle S\rangle\) for the smallest subgroup of \(G\) containing \(S\). Let \(\mathcal P_{\mathrm{prop}}(G)\) denote the proper parabolic subgroups of \(G\), relative to its Chevalley group structure.

\textbf{Classification problem (V. M. Levchuk, 1986).} For every such choice of \(G,U,V\), describe all members of
\[
\mathcal H(G;U,V)=\left\{H\leq G:\ H=\bigl\langle(H\cap U)\cup(H\cap V)\bigr\rangle\ \text{and}\ \forall P\in\mathcal P_{\mathrm{prop}}(G),\ H\nsubseteq P\right\}.
\]
Thus the subgroup must be generated by its intersections with the specified opposite unipotent subgroups and must lie in no proper parabolic subgroup.

\textbf{Unresolved scope qualifications.} The supplied record does not specify the root types or central forms covered by the term ``Chevalley group,'' whether twisted groups are included, or the rank convention if twisted groups are intended. It also does not explicitly identify \(U\) and \(V\) as the unipotent radicals of opposite Borel subgroups. If that interpretation is intended, the displayed problem is to be read with that additional specification; it is not silently imposed here. The source asks to describe all such subgroups without specifying an equivalence relation for the classification.
\subsection{Short English statement}
For each Chevalley group of rank at most \(6\) over a finite field of order at most \(9\), classify the subgroups generated by their intersections with specified opposite unipotent subgroups that are contained in no proper parabolic subgroup. The precise scope of the ambient groups and of ``opposite unipotent subgroups'' requires clarification. The registry labels this problem open; that label is supplied metadata, not an independently verified assessment.
\subsection{Sources}
\begin{itemize}
\item Local UnsolvedMath registry, permanent ID \(700637\), exact source code \texttt{catalogue entry}, title ``Kourovka Notebook Problem 10.20.'' The supplied record attributes the problem to V. M. Levchuk and dates its proposal to 1986.
\item \emph{The Kourovka Notebook (Unsolved Problems in Group Theory)}, original problem 10.20, 10th issue (1986). The supplied background identifies the reference edition as the 21st issue edition (2026), \texttt{arXiv:1401.0300v48}. The supplied collection description names E. I. Khukhro and V. D. Mazurov as editors.
\item These references and the open status are supplied source information. The cited edition and the current mathematical status have not been independently verified for this draft.
\end{itemize}
\end{document}
